% Samyama — VLDB 2026 submission (PVLDB Vol. 19)
% Target: May 1, 2026 deadline
\documentclass[sigconf, nonacm]{acmart}

%% PVLDB macros
\newcommand\vldbdoi{XX.XX/XXX.XX}
\newcommand\vldbpages{XXX-XXX}
\newcommand\vldbvolume{19}
\newcommand\vldbissue{1}
\newcommand\vldbyear{2026}
\newcommand\vldbauthors{\authors}
\newcommand\vldbtitle{\shorttitle}
\newcommand\vldbavailabilityurl{https://github.com/samyama-ai/samyama-graph}
\newcommand\vldbpagestyle{plain}

\usepackage{balance}
\usepackage{booktabs}
\usepackage{subcaption}
\usepackage{listings}
\usepackage{xcolor}
\usepackage{enumitem}
% algorithm pseudocode via plain figure

%% Cypher listing style
\lstdefinelanguage{Cypher}{
  keywords={MATCH, RETURN, WHERE, CREATE, DELETE, SET, MERGE, WITH, UNWIND,
            ORDER, BY, SKIP, LIMIT, EXPLAIN, PROFILE, CALL, YIELD, AS,
            DISTINCT, OPTIONAL, AND, OR, NOT, IN, IS, NULL, TRUE, FALSE,
            EXISTS, CASE, WHEN, THEN, ELSE, END},
  sensitive=false,
  morestring=[b]',
  morestring=[b]",
  morecomment=[l]{//},
}
\lstset{
  language=Cypher,
  basicstyle=\ttfamily\footnotesize,
  keywordstyle=\bfseries\color{blue!70!black},
  stringstyle=\color{red!70!black},
  breaklines=true,
  frame=single,
  framesep=2pt,
  xleftmargin=4pt,
  xrightmargin=4pt,
  aboveskip=6pt,
  belowskip=6pt,
}

\begin{document}

% ====================================================================
\title{Samyama: A Unified Graph-Vector Database with\\Late Materialization and Graph-Native Query Planning}

\author{Madhulatha Mandarapu}
\orcid{0009-0005-2837-6725}
\affiliation{%
  \institution{VaidhyaMegha Private Limited}
  \city{Hyderabad}
  \country{India}
}
\email{madhulatha@samyama.ai}

\author{Sandeep Kunkunuru}
\orcid{0000-0002-8886-1846}
\affiliation{%
  \institution{VaidhyaMegha Private Limited}
  \city{Hyderabad}
  \country{India}
}
\email{sandeep@samyama.ai}

% ====================================================================
% ABSTRACT
% ====================================================================
\begin{abstract}
Graph databases face a fundamental tension between flexible data modeling and query performance.
Property graph systems must support dynamic schemas, variable-length path queries, and multi-hop traversals while competing with columnar and relational engines on latency.
We present \textbf{Samyama}, a graph-vector database written in Rust that addresses this tension through three core contributions:
(1)~\emph{late materialization} for property graphs, where scan operators produce lightweight \texttt{NodeRef}/\texttt{EdgeRef} tokens instead of full property clones, achieving 4.0--4.7$\times$ speedup on multi-hop traversals;
(2)~a \emph{graph-native cost-based planner} that maintains triple-level statistics (\textsc{GraphCatalog}), enumerates candidate plans via BFS from each pattern node, and selects the cheapest using cardinality-aware cost estimation---reducing plan cost by up to 73$\times$ on real workloads; and
(3)~a \emph{two-phase bulk loading} approach with node stubs, edge stubs, and multi-segment CSR compaction that reduces per-edge memory from 709 to 52~bytes (13.6$\times$), loading a 66M-node, 1B-edge biomedical graph in 130\,GB on a \$0.67/hr spot instance.

Samyama provides ${\sim}$90\% OpenCypher coverage through 35~physical operators in a Volcano iterator model, HNSW vector indexing, and RESP/HTTP dual-protocol access.
Our evaluation on commodity hardware (Mac~Mini M4, 16\,GB) demonstrates:
255K~nodes/s ingestion; 115K~queries/s at 1M~nodes;
100\% correctness on self-validated LDBC SNB Interactive (21/21), BI (16/16), Graphalytics (12/12), and FinBench (40/40);
and sub-second query latency on a 7.9M-node biomedical knowledge graph.
\end{abstract}

\maketitle

% ====================================================================
% 1. INTRODUCTION
% ====================================================================
\section{Introduction}
\label{sec:intro}

The rise of Retrieval-Augmented Generation (RAG) and knowledge-graph-grounded AI agents has renewed interest in graph databases that combine relational structure with semantic similarity search.
Industrial deployments increasingly require a single system that handles property graph traversals, vector nearest-neighbor queries, and analytical workloads---without the ETL overhead of composing Neo4j, Pinecone, and Spark into a pipeline.

Existing graph databases face two performance bottlenecks.
First, \emph{eager materialization}: systems like Neo4j clone full node/edge records at scan time, incurring allocation overhead that dominates multi-hop traversal latency.
Second, \emph{relational planning}: most graph query planners reuse join-ordering heuristics from relational optimizers, missing graph-specific opportunities such as traversal direction reversal and adjacency-driven cardinality estimation.

We present \textbf{Samyama} (Sanskrit: the three-fold union of \emph{dh\=ara\d{n}\=a}~(focused query), \emph{dhy\=ana}~(sustained analysis), and \emph{sam\=adhi}~(unified insight)~\cite{isha_samyama}), a graph-vector database written in Rust (1,927 unit tests, 85\% code coverage) that addresses both bottlenecks. Our contributions are:

\begin{enumerate}[leftmargin=*]
  \item \textbf{Late materialization for property graphs} (\S\ref{sec:late-mat}):
    Scan operators produce \texttt{Value::NodeRef(id)} tokens instead of cloning full records.
    Properties are resolved on demand via a columnar store.
    This reduces 1-hop traversal latency by 4.0$\times$ and 2-hop by 4.7$\times$ (Table~\ref{tab:late-mat}).

  \item \textbf{Graph-native cost-based planner} (\S\ref{sec:planner}):
    A \textsc{GraphCatalog} maintains triple-level statistics (label $\times$ edge-type $\times$ label $\rightarrow$ count).
    The planner enumerates candidate plans from each pattern node, scores them using multiplicative cardinality estimation, and selects the cheapest.
    On a 3-variable triangle query over 7.9M nodes, the chosen plan is 73$\times$ cheaper than the worst candidate .

  \item \textbf{Hybrid CSR adjacency} (\S\ref{sec:csr}):
    Bulk-loaded edges are compacted into a frozen Compressed Sparse Row tier ($O(1)$ offset lookup, sorted neighbor lists),
    while runtime writes go to a mutable Vec-of-Vec buffer.
    Inter-phase compaction during streaming loads keeps memory bounded---we load 40M PubMed articles on a 128\,GB instance without OOM.
\end{enumerate}

\noindent
Additionally, Samyama provides: ${\sim}$90\% OpenCypher coverage (MATCH, CREATE, MERGE, SET, DELETE, WITH, UNWIND, UNION, EXPLAIN, PROFILE); 35 physical operators; HNSW vector indexing; RESP (Redis wire protocol) and HTTP/REST dual-protocol access; RocksDB persistence with WAL; multi-tenancy via column families; and 14 graph algorithms (PageRank, WCC, SCC, BFS, Dijkstra, LCC, CDLP, triangle counting, max flow, MST, PCA).

These contributions are evaluated on standard benchmarks (LDBC SNB, BI, Graphalytics, FinBench) and real-world workloads (7.9M-node biomedical knowledge graphs, 21K-match cricket analytics).
Samyama is open-source under Apache-2.0.\footnote{\url{https://github.com/samyama-ai/samyama-graph}}
A companion architecture book with detailed design rationale is available at \url{https://samyama-ai.github.io/samyama-graph-book/}.

% ====================================================================
% 2. RELATED WORK
% ====================================================================
\section{Related Work}
\label{sec:related}

\textbf{Graph databases.}
Neo4j~\cite{neo4j2023} is the most widely deployed property graph database, using a native graph storage engine with Cypher query support.
Its JVM-based architecture incurs garbage collection pauses (10--100\,ms) and eager materialization overhead.
Memgraph~\cite{memgraph2023} is an in-memory C++ graph database achieving sub-millisecond traversals through pre-compiled query modules (MAGE), but lacks native vector search.
Kuzu~\cite{kuzu2023} introduces structured property graphs with columnar storage and factorized processing, targeting embedded analytical workloads.
DuckPGQ~\cite{duckpgq2023} extends DuckDB with SQL/PGQ graph pattern matching over relational tables, offering graph queries without a dedicated graph engine.
TigerGraph provides a distributed graph analytics platform with GSQL, achieving high throughput through pre-compiled queries and massively parallel processing.

\textbf{Late materialization.}
The concept of deferring tuple reconstruction originates in columnar databases: C-Store~\cite{abadi2008column} and MonetDB~\cite{boncz2005monetdb} scan column projections and join position lists before materializing full rows.
We adapt this principle to property graphs, where the ``row'' is a node or edge with variable-schema properties.
To our knowledge, Samyama is the first property graph system to implement late materialization at the Volcano operator level.

\textbf{Graph query optimization.}
The GF-RBO approach in Kuzu~\cite{kuzu2023} performs factorized processing to avoid intermediate result explosion.
Worst-case optimal (WCO) join algorithms~\cite{ngo2018wco} provide theoretical guarantees for cyclic queries via LeapFrog TrieJoin.
Our planner takes a complementary approach: BFS-based plan enumeration with graph-native cardinality estimation from triple-level statistics, leveraging sorted adjacency lists for efficient \textsc{ExpandInto} operations.
GraphflowDB~\cite{graphflowdb2023} uses worst-case optimal joins with factorized processing; our system currently uses hash joins and plans to add WCO joins as future work.

\textbf{Graph-vector integration.}
Recent systems like Weaviate and Milvus add graph-like features to vector databases, while Neo4j~5.x adds vector index support via Lucene.
Samyama takes the opposite approach: a graph-first architecture with native HNSW vector indexing as a secondary capability, enabling Graph RAG patterns (vector-retrieve-then-traverse) within a single engine.

\textbf{Rust-based databases.}
The Rust programming language has gained traction in database systems for its memory safety guarantees and zero-cost abstractions.
SurrealDB provides a multi-model database in Rust, but targets document/graph hybrid workloads without a dedicated graph query optimizer.
Databend (analytics), TiKV (key-value), and Neon (Postgres) demonstrate Rust's viability for production database systems.
Samyama is, to our knowledge, the first Rust-based property graph database with a cost-based query planner and LDBC-compatible self-validation.

Table~\ref{tab:related-comparison} summarizes the feature matrix across systems.

\begin{table}[!htbp]
\caption{Feature comparison with related graph database systems.}
\label{tab:related-comparison}
\begin{tabular}{lccccc}
\toprule
\textbf{Feature} & \rotatebox{60}{\textbf{Samyama}} & \rotatebox{60}{\textbf{Neo4j}} & \rotatebox{60}{\textbf{Memgraph}} & \rotatebox{60}{\textbf{Kuzu}} & \rotatebox{60}{\textbf{DuckPGQ}} \\
\midrule
OpenCypher & ${\sim}$90\% & Full & Full & Subset & SQL/PGQ \\
Cost-based planner & \checkmark & \checkmark & \checkmark & \checkmark & \checkmark \\
Late materialization & \checkmark & \texttimes & \texttimes & \checkmark$^*$ & \checkmark$^*$ \\
Native vector search & \checkmark & v5.x & \texttimes & \texttimes & \texttimes \\
Graph algorithms & 14 & GDS & MAGE & \texttimes & \texttimes \\
LDBC-compatible & \checkmark & \checkmark & \checkmark & \checkmark & Partial \\
Zero-GC & \checkmark & \texttimes & \checkmark & \checkmark & \checkmark \\
Embedded mode & \checkmark & \texttimes & \texttimes & \checkmark & \checkmark \\
\bottomrule
\end{tabular}
\\[4pt]
{\footnotesize $^*$Columnar storage provides partial late materialization in a different form.}
\end{table}

% ====================================================================
% 3. SYSTEM ARCHITECTURE
% ====================================================================
\section{System Architecture}
\label{sec:architecture}

Samyama is a single-binary database server providing dual-protocol access: RESP (Redis wire protocol) on port 6379 for \texttt{GRAPH.QUERY} commands compatible with FalkorDB clients, and HTTP on port 8080 for REST API, Prometheus metrics (\texttt{/metrics}), health checks (\texttt{/health}), and a built-in query visualizer.
Figure~\ref{fig:architecture} shows the overall system architecture.

\begin{figure}[t]
  \centering
  \includegraphics[width=\linewidth]{figures/architecture.pdf}
  \caption{Samyama system architecture. OpenCypher queries flow through a PEG parser, cost-based planner, and Volcano-model executor to the hybrid CSR/Vec-of-Vec storage layer with RocksDB persistence.}
  \label{fig:architecture}
\end{figure}

\subsection{Storage Engine}

The storage layer uses \textbf{RocksDB}~\cite{rocksdb} for persistence with LZ4/Zstd compression and column families for tenant isolation.
In-memory, \texttt{NodeId} and \texttt{EdgeId} are direct \texttt{u64} indices into contiguous arena storage, providing $O(1)$ access.
Properties are stored in both a per-node HashMap (for point lookups) and a columnar store (for scan-heavy workloads), with automatic synchronization via \texttt{set\_node\_property()} which updates both stores atomically.

Adjacency lists use the hybrid CSR structure described in \S\ref{sec:csr}: a frozen tier for bulk-loaded data and a mutable write buffer for runtime operations.
Both outgoing and incoming adjacency lists are maintained (2$\times$ storage) for bidirectional traversal.

\subsection{Query Processing}

Queries are parsed by a Pest-based PEG grammar~\cite{ford2004peg} providing ${\sim}$90\% OpenCypher coverage, including:
MATCH (with OPTIONAL), CREATE, MERGE, SET, DELETE, REMOVE, WITH, UNWIND, UNION, RETURN~DISTINCT, ORDER~BY, SKIP, LIMIT, EXPLAIN, PROFILE, CALL procedures, CASE expressions, EXISTS subqueries, parameterized queries (\texttt{\$param}), and list/map literals.
The grammar uses atomic rules for keyword operators (e.g., \texttt{@\{  \^{}"AND" $\sim$ !(ASCII\_ALPHANUMERIC)\}}) to enforce word boundaries---a non-obvious requirement in PEG parsers where \texttt{\^{}"OR"} would otherwise match the prefix of ``ORDER''.

The executor follows the Volcano iterator model~\cite{graefe1994volcano} with 35 physical operators organized across 9 categories:
\emph{scan} (NodeScan, IndexScan),
\emph{traversal} (Expand, ExpandInto, ShortestPath),
\emph{filter} (Filter),
\emph{join} (HashJoin),
\emph{aggregation} (GroupBy, WithBarrier),
\emph{sort/limit} (Sort, Limit),
\emph{write} (Create, Delete, Set, Merge, Foreach),
\emph{index} (CreateIndex, DropIndex, ShowIndexes),
and \emph{specialized} (Unwind, Union, Project).

\textbf{Supported functions} (40+): string (toUpper, toLower, trim, replace, substring, split, reverse), numeric (abs, ceil, floor, round, sqrt, sign, log, exp, sin, cos, pi, e, haversin), aggregation (count, sum, avg, min, max, collect, percentileCont, percentileDisc, stDev), list (size, head, last, tail, range, keys, reduce), type (id, labels, type, elementId, valueType, properties, toInteger, toFloat, toBoolean), and temporal (datetime, date, time, duration with component access).

\subsection{PEG Grammar Design}

The OpenCypher grammar is implemented as a Parsing Expression Grammar (PEG) using the \texttt{pest} crate~\cite{ford2004peg}.
PEGs differ from context-free grammars in their ordered-choice operator: given alternatives $A / B / C$, the parser tries $A$ first and commits if it matches.
This determinism eliminates ambiguity but introduces two subtle challenges for Cypher:

\textbf{Keyword word boundaries.}
In PEG, \texttt{\^{}"OR"} matches the literal ``OR'' but has no word boundary check---it will match the ``OR'' prefix of ``ORDER''.
Non-atomic rules (\texttt{\{ \}}) insert implicit whitespace between elements, so \texttt{\{ \^{}"AND" $\sim$ !(ALPHA) \}} consumes the space between AND and the next word before the negative lookahead fires.
The solution: use atomic rules (\texttt{@\{ \^{}"AND" $\sim$ !(ASCII\_ALPHANUMERIC | "\_") \}}) for keyword operators, preventing implicit whitespace consumption before the boundary check.

\textbf{Operator ordering.}
In the \texttt{comparison\_op} rule, \texttt{<>} must appear before \texttt{<}, since PEG tries alternatives left-to-right.
Similarly, in the \texttt{primary} rule, \texttt{value} (which matches \texttt{TRUE}/\texttt{FALSE}/\texttt{NULL}) must appear before \texttt{variable}, otherwise these keywords are parsed as variable names.

These lessons, documented in our architecture decision records (ADR-013), apply broadly to PEG-based query language implementations.

\subsection{Vector Search}

Samyama integrates HNSW (Hierarchical Navigable Small World) vector indexing for nearest-neighbor search.
Vector indexes are created per-label and support cosine, L2, and inner-product distance metrics.
This enables Graph~RAG patterns within a single engine: retrieve similar nodes via vector search, then traverse the graph to gather structured context---without cross-system ETL.

\subsection{Graph Algorithms}

The \texttt{samyama-graph-algorithms} crate provides 14 algorithms accessible via Cypher \texttt{CALL} procedures:
PageRank (with dangling node redistribution), WCC (Union-Find), SCC (Tarjan), CDLP (label propagation), LCC (local clustering coefficient), triangle counting, BFS, Dijkstra, all-pairs shortest paths, Edmonds-Karp max flow, Prim's MST, and PCA (randomized SVD).
Algorithms operate on a CSR-projected \texttt{GraphView} for cache-efficient traversal with \texttt{rayon}-based parallelism.

\subsection{Multi-Tenancy}

Each tenant operates on an isolated \texttt{GraphStore} instance with configurable resource quotas (max nodes, max edges, max query time, max memory bytes).
Tenant data is persisted in separate RocksDB column families, enabling independent backup and recovery.
The HTTP API supports tenant-scoped queries via \texttt{/api/query} with a \texttt{graph} parameter, and tenant-scoped snapshot import/export via \texttt{/api/tenants/:id/snapshot/import}.

% ====================================================================
% 4. LATE MATERIALIZATION
% ====================================================================
\section{Late Materialization for Property Graphs}
\label{sec:late-mat}

Traditional property graph engines materialize full node records at scan time: the scan operator clones all properties into each output row.
For multi-hop traversals, intermediate hops produce records that are immediately consumed by the next expand---properties are cloned but never read.
This eager materialization is a significant overhead: for a 2-hop traversal returning 10 results from a 100K-node graph, the baseline approach clones millions of property maps that are immediately discarded.

\subsection{Design}

Samyama introduces two lightweight value types:

\begin{itemize}[leftmargin=*]
  \item \texttt{Value::NodeRef(NodeId)} --- an 8-byte token representing an unresolved node.
  \item \texttt{Value::EdgeRef(EdgeId)} --- an 8-byte token representing an unresolved edge.
\end{itemize}

\noindent
Scan operators (\texttt{NodeScanOperator}, \texttt{IndexScanOperator}) produce \texttt{NodeRef} tokens instead of full \texttt{Node} clones.
The expand operator (\texttt{ExpandOperator}) uses \texttt{get\_outgoing\_edge\_targets()} which returns \texttt{(EdgeId, NodeId, NodeId, \&EdgeType)} tuples---no edge cloning occurs.

Property access (e.g., \texttt{n.name} in a WHERE clause or RETURN) invokes \texttt{Value::resolve\_property(prop, store)}, which checks the columnar store first (cache-friendly sequential access) and falls back to the node's property HashMap.

Full materialization occurs only at the \texttt{ProjectOperator} when \texttt{RETURN n} requires the complete node.
For \texttt{RETURN n.name, n.age}, only the requested properties are resolved---the remaining properties are never touched.

\subsection{Compatibility}

To maintain correctness across operators that compare values (e.g., \texttt{JoinOperator} matching on node identity), \texttt{PartialEq} and \texttt{Hash} are defined such that \texttt{NodeRef(id) == Node(id, \_)}.
This ensures that a \texttt{NodeRef} from a scan matches a materialized \texttt{Node} from a different operator path, which is essential for the correctness of join operators that process records from multiple pipeline branches.

\subsection{Columnar Property Store}

Properties are stored in a \texttt{ColumnStore}: a per-label, per-property \texttt{HashMap<NodeId, PropertyValue>}.
When a property is set via \texttt{set\_node\_property()}, both the node's HashMap and the columnar store are updated atomically.
For scan-heavy workloads (e.g., \texttt{WHERE n.age > 30}), the columnar store enables predicate evaluation without loading the full node, further reducing memory bandwidth.

% ====================================================================
% 5. GRAPH-NATIVE QUERY PLANNER
% ====================================================================
\section{Graph-Native Cost-Based Planner}
\label{sec:planner}

Most graph database planners reuse relational join-ordering techniques, treating graph pattern matching as a series of hash joins.
We argue that graph queries have inherent structure---star patterns, chains, triangles---that a \emph{graph-native} planner can exploit directly.

\subsection{GraphCatalog: Triple-Level Statistics}
\label{sec:catalog}

The \textsc{GraphCatalog} maintains statistics at the triple level: for each pattern \texttt{(source-label, edge-type, target-label)}, it stores the edge count and average fan-out.
These statistics are incrementally maintained on every \texttt{CREATE} and \texttt{DELETE} operation, with no explicit \texttt{ANALYZE} step required.

The catalog exposes three estimation methods:

\begin{itemize}[leftmargin=*]
  \item \texttt{estimate\_label\_scan(label)}: returns the node count for a label. $O(1)$ via a label-count map.
  \item \texttt{estimate\_expand\_out(src\_label, edge\_type)}: returns the expected number of edges traversed when expanding from a source label via an edge type.
  \item \texttt{estimate\_edge\_existence(src, tgt, type)}: estimates the probability that a specific edge exists between two bound nodes, used for \textsc{ExpandInto} cost estimation.
\end{itemize}

\noindent
Additionally, per-property statistics (null fraction, distinct count, selectivity) are sampled from the first 1,000 nodes per label, enabling filter selectivity estimation.

\subsection{Plan Enumeration}
\label{sec:enumeration}

Given a MATCH pattern with $k$ pattern nodes, the planner generates up to $k$ candidate plans, each starting from a different pattern node:

\begin{enumerate}[leftmargin=*]
  \item Parse the MATCH clause into a \textsc{PatternGraph}: a set of pattern nodes (with optional label constraints) and pattern edges (with optional type constraints and direction).
  \item For each pattern node $v$, perform a BFS over the pattern graph, greedily selecting the cheapest next edge at each step based on \textsc{GraphCatalog} estimates.
  \item At each BFS step, choose the traversal direction (forward or reverse) that minimizes estimated cardinality. If the query specifies \texttt{(a)-[:KNOWS]->(b)} but scanning \texttt{b} and traversing incoming edges is cheaper, the planner reverses direction.
  \item If both endpoints of an edge are already bound, insert an \textsc{ExpandInto} operator (edge existence check via binary search, $O(\log d)$) instead of a full \textsc{Expand} (fan-out scan).
\end{enumerate}

\noindent
Each candidate plan is scored using a multiplicative cost model:
\[
  \text{cost}(P) = \sum_{i=1}^{n} \prod_{j=1}^{i} \text{selectivity}(op_j)
\]
\noindent
where $\text{selectivity}(op_j)$ is the estimated output cardinality of operator $j$.
The plan with the lowest cost is selected.
The planner evaluates up to 64 candidates (configurable via \texttt{PlannerConfig}).

\paragraph{Worked example.}
Consider the Cypher query:

\begin{lstlisting}
MATCH (d:Drug)-[:HAS_INDICATION]->(i:Indication)
      <-[:STUDIES]-(t:ClinicalTrial)
RETURN d.name, i.name, t.nct_id LIMIT 10
\end{lstlisting}

\noindent
The pattern graph has 3 nodes (\texttt{d}, \texttt{i}, \texttt{t}) and 2 edges.
The planner generates 3 candidate plans:

\begin{enumerate}[leftmargin=*]
  \item \textbf{Start from Indication} (2,844 nodes): Scan \texttt{i}, expand incoming \textsc{HAS\_INDICATION} to find \texttt{d} (fan-out: $14{,}744 / 2{,}844 \approx 5.2$), then expand incoming \textsc{STUDIES} to find \texttt{t}.
    Cost: $2{,}844 \times 5.2 \approx 14{,}744$.
  \item \textbf{Start from Drug} (19,842 nodes): Scan \texttt{d}, expand outgoing \textsc{HAS\_INDICATION} to \texttt{i}, then expand incoming \textsc{STUDIES} to \texttt{t}.
    Cost: $19{,}842 \times 13.5 \approx 266{,}926$.
  \item \textbf{Start from ClinicalTrial} (575,778 nodes): Scan \texttt{t}, expand outgoing \textsc{STUDIES} to \texttt{i}, then expand incoming \textsc{HAS\_INDICATION} to \texttt{d}.
    Cost: $575{,}778 \times 1.78 \approx 1{,}026{,}911$.
\end{enumerate}

\noindent
Plan~1 is selected (73$\times$ cheaper than Plan~3).
The key insight: Indication is the \emph{hub} label in this pattern---it has the fewest nodes but sits at the junction of both edge types.
A relational planner that considers only pairwise join cardinalities might miss this graph-structural insight.

\paragraph{Direction reversal.}
When the query specifies \texttt{(p:Protein)<-[:INTERACTS\_WITH]-(p2:Protein)}, both endpoints have the same label (37,990~nodes).
The planner evaluates both directions:
(a)~scan \texttt{p}, expand incoming \textsc{INTERACTS\_WITH} (cost: 618,863);
(b)~scan \texttt{p2}, expand outgoing \textsc{INTERACTS\_WITH} (cost: 637,837).
The 3\% cost difference arises from asymmetric edge fan-out in the directed graph.
The planner selects (a) and reverses the traversal direction at the physical level.

Figure~\ref{fig:enumerate} summarizes the plan enumeration procedure.

\begin{figure}[t]
\centering
\fbox{\parbox{0.93\linewidth}{\small
\textbf{Algorithm: Graph-Native Plan Enumeration}\\[4pt]
\textbf{Input:} Pattern graph $G_P = (V_P, E_P)$, GraphCatalog $C$\\
\textbf{Output:} Cheapest plan from up to $|V_P|$ candidates\\[4pt]
\textbf{for each} pattern node $v \in V_P$:\\
\quad $plan \leftarrow \textsc{LabelScan}(v)$; \quad $bound \leftarrow \{v\}$\\
\quad $queue \leftarrow$ edges adjacent to $v$ in $G_P$\\[2pt]
\quad \textbf{while} $queue \neq \emptyset$:\\
\quad\quad Pick cheapest edge $e = (u, w)$ from $queue$ via $C$\\
\quad\quad \textbf{if} both $u, w \in bound$: \quad\quad\quad\quad \textit{// triangle closure}\\
\quad\quad\quad $plan \leftarrow \textsc{ExpandInto}(plan, u, w, e.\text{type})$\\
\quad\quad \textbf{else}:\\
\quad\quad\quad $new \leftarrow$ unbound endpoint of $e$\\
\quad\quad\quad $dir \leftarrow \arg\min_{d \in \{fwd, rev\}} C.\text{estimate}(e, d)$\\
\quad\quad\quad $plan \leftarrow \textsc{Expand}(plan, \text{bound}, new, e.\text{type}, dir)$\\
\quad\quad\quad $bound \leftarrow bound \cup \{new\}$; add $new$'s edges to $queue$\\[2pt]
\quad $candidates.\text{append}\big(plan, \textsc{EstimateCost}(plan, C)\big)$\\[4pt]
\textbf{return} cheapest plan from $candidates$
}}
\caption{Plan enumeration via BFS from each pattern node. \textsc{ExpandInto} is used when both endpoints are bound (edge-existence filter); \textsc{Expand} fans out to discover new nodes. Direction is chosen by cost.}
\label{fig:enumerate}
\end{figure}

\subsection{Logical Plan IR and Optimization}

Plans are first expressed as a \texttt{LogicalPlanNode} IR with six variants:
\textsc{LabelScan}, \textsc{Expand}, \textsc{ExpandInto}, \textsc{Filter}, \textsc{Join}, and \textsc{CartesianProduct}.
A rule-based logical optimizer then applies three transformations:

\begin{itemize}[leftmargin=*]
  \item \textbf{Predicate pushdown}: Filter predicates referencing only one variable are pushed below Expand operators, reducing intermediate cardinality before fan-out.
  \item \textbf{ExpandInto insertion}: When both endpoints of an edge pattern are already bound by prior operators, Expand is replaced with ExpandInto---semantically a filter (does this edge exist?) rather than a fan-out.
  \item \textbf{Early LIMIT propagation}: LIMIT operators are pushed into scan and expand operators to enable short-circuit evaluation, avoiding full scans when only $k$ results are needed.
\end{itemize}

\noindent
The optimized logical plan is then converted to a physical plan via \texttt{logical\_to\_physical()}, which maps logical operators to physical Volcano-model operators.
Direction reversal is handled at this stage: a logical \texttt{Expand\{direction: Reverse\}} becomes a physical \texttt{ExpandOperator} with \texttt{Direction::Incoming}.

\subsection{Plan Cache}

Samyama caches physical plans keyed by normalized query strings with a generation-based invalidation strategy.
When the graph is modified (CREATE, DELETE, SET), the generation counter increments, invalidating all cached plans.
This provides a simple correctness guarantee without fine-grained dependency tracking, at the cost of invalidating plans that may still be valid after unrelated mutations.

\subsection{EXPLAIN and PROFILE}

\texttt{EXPLAIN} returns the physical plan as an operator tree with cost estimates, planner diagnostics (number of candidates evaluated, cost of each), and graph statistics.
\texttt{PROFILE} executes the query and annotates each operator with actual row counts and execution time, enabling developers to identify bottlenecks and validate the planner's cardinality estimates.

Figure~\ref{fig:explain} shows the \texttt{EXPLAIN} output for a 3-variable query over the biomedical dataset.
The output includes: (1)~the chosen plan as an operator tree; (2)~diagnostics showing all 3 candidates with costs; and (3)~graph statistics (19 labels, 22 edge types) used by the cost model.
The built-in Plan Visualizer in Samyama Insight renders these as interactive trees with side-by-side comparison of all candidate plans.

\begin{figure}[t]
\centering
\fbox{\parbox{0.93\linewidth}{\ttfamily\footnotesize
\textbf{EXPLAIN} MATCH (d:Drug)-[:HAS\_IND]->(i:Indication)\\
\hspace*{2em}<-[:STUDIES]-(t:ClinicalTrial)\\
\hspace*{2em}RETURN d.name, i.name, t.nct\_id LIMIT 10\\[6pt]
Limit (10)\\
+- Project (d.name, i.name, t.nct\_id)\\
\hspace*{1em}+- Expand ((i)<-[:STUDIES]-(t))\\
\hspace*{2em}+- Expand ((i)<-[:HAS\_IND]-(d))\\
\hspace*{3em}+- NodeScan (var=i, labels=["Indication"])\\[4pt]
--- Planner Diagnostics ---\\
Candidates evaluated: 3\\
Chosen plan cost: 14,744.00\\
Alternative plans:\\
\hspace*{1em}\#2 (cost 266,925.59): Drug $\to$ Indication $\to$ Trial\\
\hspace*{1em}\#3 (cost 1,026,910.74): Trial $\to$ Indication $\to$ Drug\\[4pt]
--- Statistics ---\\
Total nodes: 7,925,858 \quad Total edges: 28,000,752\\
Avg out-degree: 3.53 \quad Labels: 19 \quad Edge types: 22
}}
\caption{EXPLAIN output showing the chosen plan, 3 candidate costs, and graph statistics.
The planner starts from Indication (2.8K nodes)---73$\times$ cheaper than ClinicalTrial (575K).}
\label{fig:explain}
\end{figure}

% ====================================================================
% 6. HYBRID CSR ADJACENCY
% ====================================================================
\section{Hybrid CSR Adjacency Storage}
\label{sec:csr}

\subsection{Motivation}

The standard adjacency representation in graph databases is a Vec-of-Vec: an outer vector indexed by node ID, with each inner vector storing sorted neighbor pairs \texttt{(NodeId, EdgeId)}.
For 80M nodes $\times$ 2 directions, the inner-vector headers alone consume 3.84\,GB (24 bytes each: pointer + length + capacity), plus capacity-rounding waste from Rust's geometric growth strategy.

\begin{figure}[t]
  \centering
  \includegraphics[width=0.85\linewidth]{figures/csr_layout.pdf}
  \caption{CSR data layout. \texttt{offsets[i]} and \texttt{offsets[i+1]} delimit node~$i$'s neighbor range in the packed \texttt{edges} array. Binary search within each range provides $O(\log d)$ edge-existence checks.}
  \label{fig:csr}
\end{figure}

\subsection{FrozenAdjacency: CSR Design}

We introduce \texttt{FrozenAdjacency}, a Compressed Sparse Row (CSR) structure:

\begin{itemize}[leftmargin=*]
  \item \texttt{offsets: Vec<u32>} --- position $i$ gives the start index in \texttt{edges} for node $i$'s neighbors. Length = max\_node\_id + 2.
  \item \texttt{edges: Vec<(NodeId, EdgeId)>} --- packed, sorted neighbor entries.
\end{itemize}

\noindent
Neighbor lookup is $O(1)$ (two offset reads + slice); edge existence check is $O(\log d)$ via binary search on the sorted neighbor list.
Memory overhead per node is 4 bytes (one \texttt{u32} offset) instead of 24 bytes (one \texttt{Vec} header), a 6$\times$ reduction.

\subsection{Two-Tier Architecture}

The graph maintains both tiers simultaneously:

\begin{itemize}[leftmargin=*]
  \item \textbf{Frozen tier} (CSR): bulk-loaded data, immutable, memory-efficient.
  \item \textbf{Write buffer} (Vec-of-Vec): new edges from runtime \texttt{CREATE} operations.
\end{itemize}

\noindent
Query operators merge results from both tiers transparently.
For example, \texttt{get\_outgoing\_edges(node)} concatenates the frozen tier's slice with the write buffer's vector for that node.
\texttt{edge\_between(src, tgt)} checks the frozen tier (binary search) first, then falls back to the write buffer (linear scan).

\subsection{Compaction}

\texttt{compact\_adjacency()} freezes the write buffer into CSR:

\begin{enumerate}[leftmargin=*]
  \item If the frozen tier is empty (first compaction): convert Vec-of-Vec directly to CSR via \texttt{FrozenAdjacency::from\_vec\_of\_vec()}.
  \item If the frozen tier has data (subsequent compaction): merge frozen + buffer into a new Vec-of-Vec, sort, and re-freeze.
  \item Clear the write buffer and call \texttt{shrink\_to\_fit()} on each inner vector, releasing allocated capacity back to the system allocator.
\end{enumerate}

\noindent
For streaming bulk loads (e.g., PubMed: 40M articles, 1B edges), we call \texttt{compact\_adjacency()} between loading phases---analogous to GC pauses in managed runtimes.
This prevents the write buffer from growing unbounded, keeping peak memory within the frozen CSR footprint plus one phase's worth of buffered edges.
The trade-off is a brief compaction pause (${\sim}$10--30s for 200M edges) between phases.

\subsection{Memory Analysis}

Table~\ref{tab:memory} quantifies the memory savings from hybrid CSR for different graph scales.

The CSR representation achieves a \textbf{6$\times$ header reduction}: 4-byte \texttt{u32} offsets replace 24-byte \texttt{Vec} headers per node.
For the 66M-node PubMed graph with bidirectional adjacency, this saves $2 \times 1.6\,\text{GB} = 3.2\,\text{GB}$ of overhead.
Additionally, the Vec-of-Vec approach incurs capacity rounding waste (Rust doubles allocation on grow), adding 20--40\% overhead that CSR eliminates.

% ====================================================================
% 7. DEPLOYMENT AND DISTRIBUTION
% ====================================================================
\section{Deployment}
\label{sec:deployment}

\subsection{High Availability}

Samyama implements Raft consensus~\cite{ongaro2014raft} via the \texttt{openraft} crate for leader election and log replication.
The Enterprise Edition adds HTTP/2 Raft transport with TLS, automated snapshot streaming to lagging followers, and cluster metrics (role, term, replication lag).

We validated the distributed architecture on a 3-node AWS spot cluster (m6a.large, ap-south-1, \$2.23/day total):
4 tenants with 426K nodes and 1.39M edges.
All 52 benchmark queries returned consistent results across all 3 nodes within ${\sim}$100\,ms variance.
Single-node and dual-node failure simulations confirmed that surviving nodes remained fully operational.

A key finding: snapshot imports via HTTP API are in-memory only---not persisted through WAL/RocksDB.
Restarted nodes lose imported data unless the snapshot is re-imported.
WAL-based persistence for snapshot imports is planned.

\subsection{Operational Features}

The Enterprise Edition provides:
\emph{monitoring} (Prometheus metrics: query count, latency percentiles, active connections, storage sizes);
\emph{audit logging} (query text, user, tenant, latency for compliance);
\emph{backup/recovery} (full snapshots via RocksDB \texttt{BackupEngine}, incremental WAL-based deltas, point-in-time recovery with microsecond precision);
and \emph{license hardening} (Ed25519-signed JET tokens with machine fingerprint binding, clock drift protection, and signed revocation lists).

\subsection{SDK Ecosystem}

Samyama provides four client SDKs:

\begin{itemize}[leftmargin=*]
  \item \textbf{Rust SDK} (\texttt{samyama-sdk}): \texttt{EmbeddedClient} (in-process, zero-copy) and \texttt{RemoteClient} (HTTP). Extension traits for algorithms and vector search.
  \item \textbf{Python SDK} (PyO3): \texttt{SamyamaClient.embedded()} and \texttt{.connect(url)}. Direct algorithm methods (PageRank, WCC, BFS, Dijkstra).
  \item \textbf{TypeScript SDK}: HTTP client with \texttt{connectHttp()} factory.
  \item \textbf{CLI}: \texttt{samyama-cli} with query, status, ping, and REPL subcommands. Table/JSON/CSV output.
\end{itemize}

\noindent
The embedded client pattern enables in-process graph construction (e.g., ETL loaders, benchmarks, unit tests) without network overhead, while the remote client provides the same API over HTTP for deployed servers.

% ====================================================================
% 8. KNOWLEDGE GRAPH APPLICATIONS
% ====================================================================
\section{Knowledge Graph Applications}
\label{sec:applications}

To validate Samyama's suitability for real-world knowledge graph workloads beyond synthetic benchmarks, we constructed and evaluated six domain-specific knowledge graphs spanning healthcare, industry, sports, and infrastructure.
Table~\ref{tab:kgs} summarizes the scale and domain of each.

\begin{table}[!htbp]
\caption{Knowledge graphs evaluated on Samyama.}
\label{tab:kgs}
\begin{tabular}{lrrl}
\toprule
\textbf{Knowledge Graph} & \textbf{Nodes} & \textbf{Edges} & \textbf{Domain} \\
\midrule
Clinical Trials (AACT) & 7.7M & 27.0M & Healthcare \\
Biological Pathways & 119K & 835K & Bioinformatics \\
Drug Interactions & 33K & 192K & Pharmacology \\
Cricket (Cricsheet) & 37K & 1.39M & Sports analytics \\
AssetOpsBench & 12.6K & 12.6K & Industrial IoT \\
Disease Surveillance & 217K & 241K & Public health \\
\midrule
\textbf{Trifecta (federated)} & \textbf{7.9M} & \textbf{28.0M} & \textbf{Biomedical} \\
\bottomrule
\end{tabular}
\end{table}

The \textbf{Trifecta} configuration federates the first three KGs into a single tenant, demonstrating cross-domain query capabilities.
For example, a query tracing the path from a Drug through its Indication to related ClinicalTrials spans two independently constructed KGs:

\begin{lstlisting}
MATCH (d:Drug)-[:HAS_INDICATION]->(i:Indication)
      <-[:STUDIES]-(t:ClinicalTrial)
RETURN d.name, i.name, count(t)
ORDER BY count(t) DESC LIMIT 10
\end{lstlisting}

\noindent
This query executes in ${\sim}$10\,ms on the 7.9M-node federated graph, demonstrating that cross-KG traversals incur no penalty over single-KG queries---all data resides in a single \texttt{GraphStore} instance.

The \textbf{PubMed} graph (80M nodes, 1.08B edges) represents our scale target: the complete PubMed/MEDLINE baseline with 40M articles, 30M deduplicated authors, 342M MeSH annotations, 397M citation links, 66M chemical mentions, and 14M grant associations.
This graph, when complete, will be the largest open-source biomedical knowledge graph.

Each KG was constructed from public data sources using purpose-built Rust loaders (50--300 lines each) that use the \texttt{GraphStore} API directly, bypassing Cypher parsing for maximum ingestion throughput.

\subsection{Data Loading Patterns}

We observe three loading patterns in practice:

\textbf{Pattern 1: Rust loader (batch ETL).}
For large-scale loading (millions of nodes), Rust loaders call \texttt{graph.create\_node()}/\texttt{create\_edge()} directly on the \texttt{GraphStore}, avoiding Cypher parse overhead.
The AssetOpsBench loader (12.6K nodes) completes in 0.1\,s; the PubMed loader (40M articles) achieves 113K~nodes/s.
Deduplication uses in-memory \texttt{HashMap<String, NodeId>} for entity resolution (e.g., PMID $\to$ NodeId, author name $\to$ NodeId).

\textbf{Pattern 2: Cypher-based loading (small/medium).}
For smaller datasets or interactive use, \texttt{CREATE} and \texttt{MERGE} statements via the RESP or HTTP API provide Cypher-level flexibility at the cost of per-statement parse overhead.
The Python SDK uses this pattern for the assetops-kg ETL (8 steps, ${\sim}$12K Cypher statements), completing in ${\sim}$3 minutes---adequate for small graphs but 1800$\times$ slower than the Rust loader for the same data.

\textbf{Pattern 3: Snapshot import (instant).}
The \texttt{.sgsnap} format (portable, self-contained binary) enables sub-second graph loading via HTTP multipart upload:

\begin{lstlisting}[language=bash,keywordstyle=\color{black}]
curl -X POST /api/tenants/my-graph/snapshot/import
     -F "file=@pathways.sgsnap"
\end{lstlisting}

\noindent
The trifecta snapshot (7.9M nodes, 28M edges, ${\sim}$730\,MB) imports in ${\sim}$150\,s---including CSR compaction.
This pattern is ideal for deploying pre-built knowledge graphs to new servers or tenants.

\subsection{Cross-Domain Federation}

A key design advantage of the single-engine architecture is effortless cross-domain federation.
Loading multiple KGs into one tenant creates a \emph{federated} graph where queries can traverse edges across domains without any federation middleware:

\begin{lstlisting}
-- Cross-KG: Drug (DrugInteractions KG) -> Indication
--           -> ClinicalTrial (ClinicalTrials KG)
MATCH (d:Drug)-[:HAS_INDICATION]->(i)
      <-[:STUDIES]-(t:ClinicalTrial)
RETURN d.name, t.nct_id LIMIT 10
\end{lstlisting}

\noindent
The planner treats federated data identically to single-source data---\textsc{GraphCatalog} statistics are maintained across all labels regardless of their ETL origin.
This contrasts with federated query systems that require explicit routing and reconciliation.

Our BiomedQA benchmark~\cite{mandarapu2025biomedkg} evaluates 40 questions across three federated biomedical KGs, achieving 98\% accuracy with MCP-based deterministic handlers and 85\% with NLQ (text-to-Cypher).

% ====================================================================
% 9. EVALUATION
% ====================================================================
\section{Evaluation}
\label{sec:eval}

We evaluate Samyama on commodity hardware: Mac~Mini~M4 (Apple~M4, 10-core, 16\,GB LPDDR5X, macOS~Tahoe~26.2) for primary benchmarks, and AWS r6a.2xlarge (8~vCPU AMD~EPYC, 62\,GB RAM) for scale tests.
Each measurement is the median of 5 runs unless otherwise noted.

\subsection{LDBC Self-Validation}
\label{sec:ldbc}

\begin{table}[!htbp]
\caption{LDBC benchmark validation. All tests pass (100\%).}
\label{tab:ldbc}
\begin{tabular}{lrrl}
\toprule
\textbf{Benchmark} & \textbf{Pass} & \textbf{Total} & \textbf{Dataset} \\
\midrule
SNB Interactive & 21 & 21 & SF1 (3.18M nodes, 17.3M edges) \\
SNB Business Intelligence & 16 & 16 & SF1 (4 timeout on heavy global) \\
Graphalytics & 12 & 12 & XS (2 datasets) + S (3 datasets) \\
FinBench & 40 & 40 & Synthetic (7.7K nodes, 42.2K edges) \\
\bottomrule
\end{tabular}
\end{table}

Table~\ref{tab:ldbc} shows 100\% correctness across all four LDBC benchmark suites (self-validated against reference outputs).
Graphalytics S-size datasets include cit-Patents (3.8M vertices, 16.5M edges), datagen-7\_5-fb (633K vertices, 68.4M edges), and wiki-Talk (2.4M vertices, 5.0M edges).
LDBC SNB BI runs all 20 queries; BI-17 through BI-20 involve global analytics over the full graph that time out at 60s, but produce correct partial results.

\textbf{LDBC SNB Interactive} covers 21 read queries: 7 short lookups (IS1--IS7) and 14 complex reads (IC1--IC14).
Complex reads test multi-hop traversals with aggregation---for example, IC5 (``new groups of a person's friends'') requires a 2-hop traversal through KNOWS edges followed by group membership lookup and temporal filtering.
All 21 queries produce correct results on the SF1 dataset (3.18M nodes, 17.3M edges representing a social network).

\textbf{LDBC FinBench} models financial transaction networks with accounts, transfers, loans, companies, and persons.
All 40 queries (12 complex reads CR1--CR12, 6 simple reads SR1--SR6, 3 read-write queries RW1--RW3, and 19 write queries W1--W19) pass on synthetic SF1 data.
The benchmark validates temporal range queries, cyclical transaction detection, and multi-entity relationship patterns.

\textbf{LDBC Graphalytics} validates 6 graph algorithms (BFS, PageRank, WCC, CDLP, LCC, SSSP) against reference outputs.
A critical implementation detail: LDBC reference PageRank uses exact iteration counts (no early termination) with dangling-node mass redistribution.
Running extra iterations changes absolute scores even with redistribution enabled, requiring precise alignment with the reference configuration.

\subsection{Late Materialization Impact}
\label{sec:late-mat-eval}

\begin{table}[!htbp]
\caption{Late materialization impact and ablation (100K nodes, avg degree~10).}
\label{tab:late-mat}
\begin{tabular}{lrrl}
\toprule
\textbf{Configuration} & \textbf{1-hop} & \textbf{2-hop} & \textbf{Contribution} \\
\midrule
Baseline (full clone) & 164\,ms & 1,220\,ms & --- \\
+ NodeRef (lazy) & ${\sim}$60\,ms & ${\sim}$350\,ms & 2.7$\times$ (allocation) \\
+ ColumnStore & 41\,ms & 259\,ms & 1.5$\times$ (cache) \\
\midrule
\textbf{Total speedup} & \textbf{4.0$\times$} & \textbf{4.7$\times$} & \\
\bottomrule
\end{tabular}
\end{table}

Table~\ref{tab:late-mat} shows two complementary optimizations:
(1)~replacing full node clones with \texttt{NodeRef} references eliminates allocation and memcpy overhead (2.7$\times$);
(2)~the columnar property store enables cache-friendly sequential access during property resolution (additional 1.5$\times$).
The speedup grows with hop count because deeper traversals produce exponentially more intermediate records that are cloned-then-discarded in the baseline.

\subsection{Planner Effectiveness}
\label{sec:planner-eval}

We evaluated planner effectiveness on three query patterns over the biomedical trifecta dataset (7.9M nodes, 28M edges).
For the 3-variable join (Drug--Indication--ClinicalTrial), the planner evaluates 3 candidates: starting from Indication (2,844 nodes, cost~14,744), Drug (19.8K, cost~266,926), or ClinicalTrial (575K, cost~1,026,911)---selecting the Indication-first plan at \textbf{73$\times$ cheaper} than the worst.
For a predicate-pushdown query (\texttt{WHERE d.name = `Aspirin'}), starting from Gene (4.2K, cost~19,017) is 3$\times$ cheaper than Drug (19.8K, cost~57,127).

\subsection{Ingestion and Query Throughput}

On Mac Mini M4 (16\,GB): node ingestion at 255K nodes/s, edge ingestion at 4.2M edges/s, Cypher OLTP at 115K--116K queries/s (cached plan, 100K--1M nodes), vector search at 15,872 QPS (549\,$\mu$s, k=10, 128-dim).
Index-driven $O(1)$/$O(\log n)$ access ensures near-constant throughput as graph size increases.

\subsection{Scale: Biomedical Knowledge Graphs}
\label{sec:scale}

We evaluate scalability using two biomedical datasets:

\textbf{Trifecta} (7.9M nodes, 28M edges): Three federated knowledge graphs---Clinical Trials (7.7M nodes, 27M edges from 575K trials), Drug Interactions (33K nodes, 192K edges), and Biological Pathways (119K nodes, 835K edges)---loaded into a single tenant via snapshot import.
Query latency remains sub-second for 2-hop traversals: a \texttt{MATCH (d:Drug)-[:HAS\_SIDE\_EFFECT]->(s:SideEffect)} aggregation with \texttt{ORDER BY} and \texttt{LIMIT 10} completes in 91\,ms over 139K edges.

\textbf{PubMed} (80M nodes, 1.08B edges): Full PubMed/MEDLINE 2026 baseline---40M articles, 30M authors, 342M MeSH annotations, 397M citations.
Initial loading attempts OOM'd at 256\,GB and 384\,GB due to the per-node \texttt{HashMap} overhead: empirical measurement showed \textbf{3,785~bytes per node}---6.3$\times$ the theoretical minimum of 135~bytes---dominated by HashMap control structures, String header overhead, and allocator fragmentation.

We solved this via \textbf{two-phase bulk loading} with three techniques:
(a)~\texttt{create\_node\_stub()} allocates a lightweight node (label + identity only, no property HashMap---measured at 777~bytes/node vs.\ 3,785~bytes with HashMap);
(b)~\texttt{create\_edge\_stub()} stores only adjacency entries, skipping the full Edge object---measured at 52~bytes/edge vs.\ 709~bytes (a \textbf{13.6$\times$ reduction});
(c)~properties are written directly to the sparse \texttt{ColumnStore} via \texttt{set\_column\_property()}.

With these optimizations, the full PubMed graph (\textbf{66.2M nodes, 1.04B edges}) loaded in \textbf{3~hours~46~minutes} on an r6a.8xlarge (256\,GB, \$0.67/hr spot) with peak RSS of ${\sim}$130\,GB---well within the 245\,GB available.
Total cost: \textbf{\$2.50}.
Mid-phase compaction every 50M edges (19~CSR segments) kept the write buffer bounded.

Additionally, we linked 1,018,483 PubMed articles to ClinicalTrials.gov entries via AACT study references, creating cross-KG bridge properties that enable federated queries spanning both corpora.

Queries work unchanged because \texttt{resolve\_property()} (late materialization, \S\ref{sec:late-mat}) already checks the ColumnStore first.

Table~\ref{tab:pubmed-phases} shows the per-phase loading profile.
The key observation: article loading (Phase~1) slows from 220K/s to 113K/s as the graph grows, due to HashMap resizing for the 40M-entry PMID lookup table.
After Phase~1 compaction, edge-loading phases (2--6) maintain higher throughput since they avoid node dedup and operate primarily on adjacency list insertions.

\begin{table}[h]
\caption{PubMed loading: per-unit memory cost (measured locally at 1M scale, validated at 66M on AWS).}
\label{tab:pubmed-phases}
\begin{tabular}{lrrr}
\toprule
\textbf{Component} & \textbf{Full} & \textbf{Stub} & \textbf{Reduction} \\
\midrule
Per node & 3,785\,B & 777\,B & 4.9$\times$ \\
Per edge & 709\,B & 52\,B & 13.6$\times$ \\
\midrule
66M nodes & 237\,GB & 49\,GB & \\
1B edges & 709\,GB & 52\,GB & \\
ColumnStore + maps & --- & 25\,GB & \\
CSR (19 segments) & --- & 10\,GB & \\
\midrule
\textbf{Peak RSS} & \textbf{OOM at 512\,GB} & \textbf{${\sim}$130\,GB} & \\
\textbf{Instance cost} & \$1.01/hr & \textbf{\$0.67/hr} & \\
\textbf{Total cost} & Failed & \textbf{\$2.50} & \\
\bottomrule
\end{tabular}
\end{table}

\subsection{Comparative Evaluation}
\label{sec:comparison}

\begin{table}[!htbp]
\caption{Comparison with other graph databases.
Neo4j and Memgraph numbers are from published benchmarks; Samyama is measured on Mac~Mini~M4.}
\label{tab:comparison}
\begin{tabular}{lrrr}
\toprule
\textbf{Metric} & \textbf{Samyama} & \textbf{Neo4j 5.x} & \textbf{Memgraph 2.x} \\
\midrule
Node ingestion & 255K/s & ${\sim}$26K/s & ${\sim}$295K/s \\
1-hop traversal & 41\,ms & ${\sim}$28\,ms & ${\sim}$1.1\,ms \\
Memory (1M nodes) & 450\,MB & ${\sim}$1,200\,MB & ${\sim}$600\,MB \\
GC pauses & 0\,ms & 10--100\,ms & 0\,ms \\
Native vector search & \checkmark & \texttimes\ (Lucene) & \texttimes\ (ext.) \\
\bottomrule
\end{tabular}
\end{table}

\textbf{Caveats:} Direct comparison is approximate due to different hardware and datasets.
Memgraph's 1.1\,ms 1-hop reflects pre-compiled query modules; Samyama's 41\,ms includes parse~(54\%) + plan~(44\%) overhead, which the AST and plan caches reduce for repeated queries.
The raw storage layer achieves 15\,$\mu$s for 3-hop traversals, confirming the bottleneck is in query compilation.

Samyama's advantages lie in native vector search (549\,$\mu$s, unavailable natively in Neo4j/Memgraph), lower memory footprint (450\,MB vs.~1,200\,MB for Neo4j at 1M nodes due to arena allocation and zero GC overhead), and integrated graph algorithms (14 algorithms accessible via Cypher \texttt{CALL} procedures).

\textbf{Why is Samyama's 1-hop slower than Neo4j?}
Neo4j's query compiler has had 15+ years of optimization.
Samyama's PEG parser, while correct and maintainable, generates a fresh parse tree on each query (54\% of latency).
The cost-based planner adds another 44\%.
The raw storage traversal (2\% of latency) is competitive.
With AST cache (eliminates parse) and plan cache (eliminates planning), repeated queries achieve ${\sim}$15\,$\mu$s---faster than Neo4j's reported 28\,ms because the storage layer bypasses Java object creation and GC overhead.

\textbf{Why is Samyama's 1-hop slower than Memgraph?}
Memgraph's MAGE pre-compiled query modules bypass parsing entirely.
Samyama's equivalent (plan cache + AST cache for repeated queries) achieves similar benefits but requires the first execution to pay full parse+plan cost.
A future pre-compilation API (analogous to SQL prepared statements) would close this gap for known query patterns.

\textbf{Memory advantage.}
Samyama's 450\,MB for 1M nodes (vs.\ Neo4j's 1,200\,MB) comes from three factors:
(1)~arena allocation eliminates per-object memory headers and GC metadata;
(2)~\texttt{NodeId}/\texttt{EdgeId} as direct \texttt{u64} indices avoid pointer-based storage;
(3)~hybrid CSR eliminates Vec headers for frozen adjacency data.
This is a structural advantage of Rust's ownership model---memory layout is controlled at compile time.

\subsection{Real-World Workload: AssetOpsBench}
\label{sec:assetops}

To evaluate Samyama's effectiveness as a data layer for LLM-based agents, we benchmark against IBM's AssetOpsBench~\cite{patel2025assetops}---139 industrial maintenance scenarios covering IoT telemetry, failure mode reasoning, work order management, and multi-agent coordination.

\begin{table}[!htbp]
\caption{AssetOpsBench results: KG vs. flat document stores.}
\label{tab:assetops}
\begin{tabular}{lrrr}
\toprule
\textbf{Approach} & \textbf{Pass Rate} & \textbf{Avg Score} & \textbf{Latency} \\
\midrule
GPT-4 + flat docs (IBM) & 65\% (91/139) & --- & ${\sim}$11s \\
GPT-4 + graph NLQ & 82\% (114/139) & 0.790 & ${\sim}$5.8s \\
Deterministic graph & \textbf{99\% (137/139)} & \textbf{0.889} & \textbf{63\,ms} \\
\bottomrule
\end{tabular}
\end{table}

\noindent
The same-model comparison (GPT-4 in both cases) isolates the data model's contribution: +17pp from replacing flat document stores with a knowledge graph.
The deterministic graph handlers achieve 99\% at 63\,ms average latency and zero token cost.
On the expanded HuggingFace suite (467 scenarios across 6 domains), deterministic handlers achieve 100\% pass rate (avg score 0.848).

This result demonstrates that for structured operational domains, the \emph{data model}---not the LLM---is the primary bottleneck for agent performance.
Knowledge graphs replace LLM reasoning (error-prone counting, traversal, relationship correlation) with deterministic graph operations (Cypher queries, BFS traversals, PageRank).

\textbf{MCP integration.}
The deterministic graph handlers are exposed as MCP (Model Context Protocol) tools, enabling Claude, GPT-4, and other LLM agents to call graph operations as structured tool invocations.
The MCP server provides 9 tools: asset hierarchy queries (3), failure mode analysis (3), dependency impact analysis (2), and analytics (1, PageRank-based criticality ranking).
This architecture separates the LLM's role (intent understanding, answer synthesis) from the data layer's role (deterministic execution), yielding both higher accuracy and lower cost.

\subsection{Real-World Workload: Cricket Analytics}
\label{sec:cricket}

We evaluate Samyama on a ball-by-ball cricket knowledge graph: 21,324 matches, 12,933 players, 36,619 nodes, 1.39M edges from Cricsheet data (2004--2025).
A benchmark of 100 analytical queries spanning 5 complexity levels (L1: single-label scan through L5: multi-hop aggregation with temporal predicates) achieves 87/100 on a fresh server load.

Representative latencies: L1 (single scan) 12\,ms, L2 (aggregation) 89\,ms, L3 (2-hop traversal) 210\,ms, L4 (2-hop + DISTINCT) 450\,ms, L5 (3-hop + ORDER) 890\,ms.
L1--L4 complete under 500\,ms; L5 approaches 1 second.
13 failing queries involve sequential state dependencies (running totals) requiring session-level state---a data loading issue, not a query engine limitation.

\subsection{Discussion}
\label{sec:discussion}

\textbf{Late materialization trade-offs.}
The 4.0--4.7$\times$ speedup comes at a semantic cost: operators that inspect node properties (e.g., \texttt{WHERE n.name = 'Alice'}) must invoke \texttt{resolve\_property()}, adding a function call per property access.
For queries that access many properties on every node (e.g., \texttt{RETURN n} on a full scan), the lazy resolution may perform slightly worse than eager materialization due to indirection overhead.
In practice, most queries filter on 1--2 properties and return a subset, making late materialization a net win.

\textbf{Planner accuracy.}
The multiplicative cost model assumes independence between predicates and uniform value distribution.
On the biomedical dataset, we observe that the chosen plan is within 2$\times$ of the true optimal (measured via PROFILE) for 90\% of queries.
Histograms and multi-column statistics would improve accuracy for correlated predicates (e.g., \texttt{WHERE n.country = 'US' AND n.state = 'CA'}).

\textbf{CSR compaction overhead.}
The compaction pause during PubMed loading (40M nodes, 40M edges $\rightarrow$ CSR) takes 15--30 seconds per phase.
For a 6-phase load with ${\sim}$1B total edges, this adds ${\sim}$2 minutes of compaction time to a multi-hour load---acceptable for batch loading, but not suitable for high-throughput online ingestion.
The write buffer handles runtime writes without compaction pauses.

\textbf{Columnar-first as a general principle.}
The 39$\times$ per-node memory reduction from columnar-first loading is not PubMed-specific.
Any property graph with millions of nodes and a few properties per node will exhibit the same HashMap overhead.
The key insight: for in-memory property graphs, the per-node \texttt{HashMap} is an anti-pattern at scale.
Systems like Kuzu and DuckDB avoid this by using columnar storage exclusively;
our contribution is showing that a property graph system can \emph{incrementally adopt} columnar-first loading
while retaining the traditional HashMap path for small graphs and interactive use.

\textbf{Parse/plan bottleneck.}
Profiling reveals that PEG parsing accounts for 54\% and planning for 44\% of Cypher query latency, with execution at only 2\%.
The AST cache (v0.5.10) eliminates re-parsing for repeated queries; the plan cache (v0.6.0) eliminates re-planning.
For cached queries, end-to-end latency drops to ${\sim}$15\,$\mu$s (raw storage traversal), comparable to pre-compiled systems.
First-execution latency remains an area for improvement.

% ====================================================================
% 8. CONCLUSION
% ====================================================================
\section{Conclusion}
\label{sec:conclusion}

We presented Samyama, a unified graph-vector database that introduces late materialization for property graphs, a graph-native cost-based planner with triple-level statistics, and two-phase bulk loading with edge stubs.
These techniques enable competitive performance on commodity hardware: 4.0--4.7$\times$ traversal speedup from late materialization, 73$\times$ plan cost reduction from graph-native planning, and a 66M-node, 1B-edge biomedical knowledge graph loaded in 130\,GB at \$2.50 total cost via 13.6$\times$ edge memory reduction.
Samyama achieves 100\% pass rates on all four LDBC benchmark suites and serves sub-second queries on a 7.9M-node biomedical knowledge graph.

\textbf{Limitations.}
Cypher parse and plan overhead (54\% + 44\% of query latency) remains significant for point queries; pre-compiled query modules (as in Memgraph) would reduce this.
The cost model uses multiplicative cardinality estimates without histogram-based selectivity or predicate correlation.
The hybrid CSR merge path rebuilds the full CSR on subsequent compactions; an in-place merge would reduce pause time.

\textbf{Future work.}
Several directions remain:
\emph{(1)~Worst-case optimal joins.}
For cyclic patterns (triangles, cliques), LeapFrog TrieJoin over sorted adjacency lists would provide worst-case optimal guarantees.
Our sorted neighbor lists (a prerequisite for WCO joins) are already in place via the CSR frozen tier.
\emph{(2)~Parallel query execution.}
Morsel-driven parallelism within a single query would leverage multi-core CPUs for scan-heavy workloads.
The current single-threaded Volcano model leaves performance on the table for queries that scan millions of nodes.
\emph{(3)~Histogram-based cardinality estimation.}
The current cost model assumes uniform value distribution.
Equi-depth histograms per property would improve selectivity estimation for skewed data distributions.
\emph{(4)~Distributed evaluation at scale.}
LDBC SNB at SF10+ (32M nodes, 176M edges) and Jepsen-style consistency testing for the Raft-based HA layer would validate production readiness.
\emph{(5)~Adaptive query execution.}
Re-optimizing mid-execution based on actual cardinalities (when estimated cardinalities diverge by $>10\times$ from actual) would help for queries where the cost model's independence assumptions break down.

\textbf{Availability.}
Samyama is open-source under Apache-2.0 at \url{https://github.com/samyama-ai/samyama-graph}.
An architecture book with design rationale and ADRs is at \url{https://samyama-ai.github.io/samyama-graph-book/}.
Benchmark scripts, data loaders, and reproduction instructions are included in the repository.

% ====================================================================
% REFERENCES
% ====================================================================
\balance
\bibliographystyle{ACM-Reference-Format}
\bibliography{samyama-vldb}

\end{document}
